\begin{abstract}
On result day, lakhs of students try to open an examination board's website at
the same moment, and these websites often fail. Boards need evidence to decide
how to build them. We present ResultBench, an open-source testbed that compares
designs of a result website under identical conditions. It runs eleven designs
on the same two processor threads: a website that builds every page from a
database, two mirror copies of it, four variants of a cache in front of the
database, results prepared in advance in an in-memory store or as one file per
student, and virtual waiting rooms in front of these. A
load generator replays a result-day surge of synthetic students who arrive
before and after the announced time, wait a limited time for an answer and
press refresh when it is slow. Every answer is checked against the true
result, so wrong, outdated and missing results are counted alongside speed and
processor use, and every configuration is repeated three times. On one
machine, the database design served only 30.8 percent of students within five
seconds at twice the load it was sized for, while results prepared in advance
served every student at ten times that load, using less than a twelfth of the
processor time per student. No design returned wrong results, a cache helped
only when filled before release, and a waiting room could serve almost every
student, but only after waits of minutes. We turn these measurements into a
decision guide for examination boards.
\end{abstract}
